\chapter{Estado del arte y fundamentación}\label{cap:estado_del_arte_y_fundamentacion}

El éxito de una plataforma tecnológica no solo radica en su correcto funcionamiento, sino en su capacidad para aportar valor real frente a las soluciones existentes y sustentarse sobre decisiones de ingeniería justificadas. Antes de iniciar el diseño arquitectónico, es clave analizar el \textbf{ecosistema competitivo} y evaluar críticamente las \textbf{herramientas de desarrollo} disponibles en la industria actual.

A lo largo de las siguientes secciones se realiza un \textbf{análisis del mercado} de plataformas de evaluación mediante cuestionarios para identificar sus limitaciones y fijar los \textbf{factores diferenciadores} de \textit{Quizzie}. Asimismo, se detalla la \textbf{fundamentación tecnológica} del proyecto, examinando las alternativas consideradas y justificando la selección del \textbf{stack definitivo}, para concluir con la síntesis estratégica del estudio previo.
\section{Análisis del mercado}\label{sec:analisis_del_mercado}

Para fundamentar la necesidad de desarrollar una nueva plataforma de evaluación, es imprescindible analizar el ecosistema de soluciones empleadas actualmente en el ámbito educativo. Aunque el mercado cuenta con soluciones consolidadas, el estudio de sus características revela limitaciones críticas en el control de la presencialidad, la exportación de datos y la asistencia al profesorado. A continuación, se examinan las tres alternativas más representativas:

\begin{itemize}
    \item \textbf{Kahoot!:} Pionera en el ámbito de la gamificación educativa, destaca por su alta capacidad de fidelización y su interfaz visual atractiva. Sin embargo, su enfoque lúdico penaliza la rigurosidad académica: prioriza la velocidad de respuesta sobre la reflexión profunda, apenas ofrece analíticas avanzadas exportables en sus versiones gratuitas y carece de cualquier mecanismo que impida a un alumno participar de forma remota compartiendo el código PIN de la sala.
    \item \textbf{Wooclap:} Orientada a un entorno universitario y corporativo más formal, ofrece una excelente variedad de tipos de preguntas y herramientas de interacción en vivo. Su principal barrera reside en la usabilidad y la agilidad de los flujos de trabajo; el proceso de diseño de cuestionarios extensos resulta monótono y manual para el docente, y no tiene funcionalidades inteligentes de automatización que agilicen la generación de contenido académico.
    \item \textbf{Cuestionarios de Enseñanza Virtual:} Las plataformas institucionales integradas ofrecen una robustez indiscutible a nivel de seguridad, evaluación formal y exportación nativa de calificaciones. No obstante, sacrifican por completo la inmediatez y la experiencia de usuario. Su interfaz rígida, la complejidad técnica para configurar sesiones en tiempo real y la imposibilidad de mitigar de forma sencilla el fraude por suplantación digital a distancia limitan su efectividad como dinámicas de aula ágiles.
\end{itemize}

Frente a este escenario, se ha sintetizado un análisis en la siguiente tabla comparativa del mercado, evaluando las características esenciales para la docencia presencial moderna. Las carencias detectadas en las alternativas comerciales e institucionales delimitan el espacio de oportunidad para una solución integral.

\begin{table}[H]
\centering
\small
\renewcommand{\arraystretch}{1.5}
\makebox[\textwidth][c]{
\begin{tabular}{l c c c c}
\bighline
\textbf{Característica} & \textbf{Kahoot!} & \textbf{Wooclap} & \textbf{Enseñanza Virtual} & \textbf{\textit{Quizzie}} \\ \bighline
Evaluación en tiempo real & \textbf{Sí} & \textbf{Sí} & Limitado & \textbf{Sí} \\ \hline
Acceso rápido (sin registro de alumnos) & \textbf{Sí} & \textbf{Sí} & No & \textbf{Sí} \\ \hline
Verificación de presencialidad física & No & No & No & \textbf{Sí} \\ \hline
Generación de test asistida por IA & No & Limitado & No & \textbf{Sí} \\ \hline
Creación ágil de cuestionarios & \textbf{Sí} & Limitado & No & \textbf{Sí} \\ \hline
Visualización de estadísticas y analíticas & \textbf{Sí} & \textbf{Sí} & Limitado & \textbf{Sí} \\ \hline
Exportación de resultados & De pago & Limitado & \textbf{Sí} & \textbf{Sí} \\ \bighline
\end{tabular}
}
\caption{Tabla comparativa del mercado.}
\label{tab:comparativa_mercado}
\end{table}

\section{Factores diferenciadores}\label{sec:factores_diferenciadores}

Como respuesta a las debilidades del mercado identificadas en la sección anterior, \textit{Quizzie} se diseña no como una alternativa lúdica adicional, sino como un \textbf{entorno equilibrado} que fusiona el \textbf{dinamismo interactivo} con la \textbf{formalidad académica}. Sus elementos innovadores y ventajas competitivas se estructuran en los siguientes pilares:

\begin{itemize}
    \item \textbf{Validación de presencia física mediante QR:} Para evitar el fraude por suplantación o que los alumnos participen a distancia, la plataforma genera un código QR único en el dispositivo del estudiante justo al finalizar el test, el cual incluye su nombre y la nota obtenida. El docente es quien escanea este código desde su propio dispositivo al terminar la clase, confirmando de manera inequívoca la asistencia física del alumno y registrando su calificación en el acto.    
    \item \textbf{Creación inteligente asistida por IA:} Para resolver el flujo de trabajo monótono de diseñar cuestionarios desde cero, \textit{Quizzie} integra un motor de Inteligencia Artificial que asiste al docente generando automáticamente baterías de preguntas coherentes a partir del temario o texto introducido, agilizando drásticamente la preparación de la clase.
    \item \textbf{Exportación directa de calificaciones:} La plataforma simplifica la gestión de las calificaciones permitiendo descargar las notas de los alumnos directamente en formatos estándar como CSV o PDF, listos para su revisión o publicación en canales oficiales.
    \item \textbf{Accesibilidad y agilidad en el aula:} Los alumnos acceden de forma inmediata sin necesidad de registros previos ni cuentas obligatorias, mientras que el profesor dispone de un panel simplificado para lanzar y controlar el flujo del cuestionario con un solo clic.
\end{itemize}

\section{Fundamentación tecnológica}\label{sec:fundamentacion_tecnologica}

A continuación, se presenta la arquitectura tecnológica que da soporte a \textit{Quizzie}. Como vista general de la infraestructura, la siguiente tabla muestra la distribución definitiva de las herramientas seleccionadas siguiendo un orden ascendente desde la persistencia hasta el despliegue del cliente.

\begin{table}[H]
\centering
\small
\renewcommand{\arraystretch}{1.5}
\setlength{\tabcolsep}{15pt}
\makebox[\textwidth][c]{%
\begin{tabular}{lp{4cm}}
\bighline
\multicolumn{2}{c}{\textbf{\textit{Stack} tecnológico seleccionado}} \\ \bighline
\textbf{Base de datos} & SQLite \\
\hline
\textbf{ORM / Validación} & SQLModel \\
\hline
\textbf{Backend} & FastAPI (Python) \\
\hline
\textbf{Lenguaje frontend} & TypeScript \\
\hline
\textbf{Frontend (SPA)} & React + Vite \\
\hline
\textbf{Estilos frontend} & CSS \textit{vanilla} \\
\hline
\textbf{Alojamiento servidor} & Render \\
\hline
\textbf{Alojamiento cliente} & Vercel \\
\hline
\textbf{Flujo de CI/CD} & GitHub Actions \\
\bighline
\end{tabular}%
}
\caption{Resumen del \textit{stack} tecnológico de Quizzie.}
\label{tab:resumen_stack}
\end{table}

\subsection{Alternativas consideradas y justificación de las decisiones}\label{subsec:alternativas_y_justificacion}

Para fundamentar la elección de cada componente del sistema, se analizan a continuación las matrices comparativas correspondientes a cada uno de los bloques lógicos de la aplicación.

\subsubsection{Persistencia de datos}
La gestión del almacenamiento del proyecto requiere evaluar la estructura de los datos frente a la complejidad de la infraestructura en la nube.

\begin{table}[H]
\centering
\small
\renewcommand{\arraystretch}{1.4}
\makebox[\textwidth][c]{%
\begin{tabularx}{1\textwidth}{l Y Y c}
\bighline
\textbf{Tecnología} & \textbf{Ventajas} & \textbf{Desventajas} & \textbf{Elección} \\ \bighline
\textbf{SQLite} & Portabilidad en un solo archivo; cero administración de red y coste económico nulo. & Bloqueo del archivo completo durante la concurrencia de escritura masiva. & \covCheck \\ \hline
PostgreSQL & Escalabilidad masiva y gestión avanzada de accesos y concurrencia. & Elevada complejidad de despliegue y administración de accesos en producción. & \covCross \\ \hline
MySQL & Motor relacional maduro, robusto y con gran soporte de la comunidad. & Costes económicos asociados a instancias en la nube innecesarios en esta fase. & \covCross \\ \hline
MongoDB & Flexibilidad de esquemas basada en documentos NoSQL y alta velocidad de lectura. & Descartado porque el dominio de datos del proyecto es puramente relacional. & \covCross \\ \bighline
\end{tabularx}%
}
\caption{Alternativas tecnológicas para la base de datos.}
\label{tab:comp_base_datos}
\end{table}

\subsubsection{Mapeo objeto-relacional (ORM)}
Se evalúa la forma de interactuar con el motor relacional buscando la mayor limpieza y eficiencia en el código.

\begin{table}[H]
\centering
\small
\renewcommand{\arraystretch}{1.5}
\makebox[\textwidth][c]{%
\begin{tabularx}{\textwidth}{l Y Y c}
\bighline
\textbf{Tecnología} & \textbf{Ventajas} & \textbf{Desventajas} & \textbf{Elección} \\ \bighline
\textbf{SQLModel} & Cumple el principio \textit{DRY}: fusiona eficientemente esquemas de API y modelos de BD. & Herramienta más reciente en el ecosistema, con menor volumen de hilos de soporte. & \covCheck \\ \hline
SQLAlchemy + Pydantic & Estándar clásico de la industria con máxima madurez y documentación. & Duplicidad de código obligatoria al tener que declarar los modelos por separado. & \covCross \\ \bighline
\end{tabularx}%
}
\caption{Alternativas tecnológicas para el ORM y validación.}
\label{tab:comp_orm}
\end{table}

\subsubsection{Backend y servidor de sockets}
La lógica de negocio exige un entorno que soporte una alta concurrencia interactiva en tiempo real.

\begin{table}[H]
\centering
\small
\renewcommand{\arraystretch}{1.5}
\makebox[\textwidth][c]{%
\begin{tabularx}{\textwidth}{l P{6.3cm} Y c}
\bighline
\textbf{Tecnología} & \textbf{Ventajas} & \textbf{Desventajas} & \textbf{Elección} \\ \bighline
\textbf{FastAPI} & Rendimiento asíncrono nativo (\textit{async/await}) según los principios de aplicaciones web modernas de Fox y Patterson \cite{fox2014engineering}, ideal para mantener conexiones \textit{WebSocket} concurrentes. & Estructura menos guiada en comparación con \textit{frameworks} monolíticos. & \covCheck \\ \hline
Node.js & Estándar tradicional para \textit{WebSockets}; unifica el lenguaje del proyecto a JavaScript. & Se descartó para priorizar el uso de Python debido a su ecosistema de IA. & \covCross \\ \hline
Django & \textit{Framework} muy robusto, maduro y con arquitectura \textit{batteries-included} basada en patrones empresariales de Fowler \cite{fowler2002patterns}. & Estructura fundamentalmente síncrona por defecto y peso excesivo para una API REST. & \covCross \\ \hline
Flask & \textit{Microframework} extremadamente simple, minimalista y rápido de configurar. & Carece de herramientas nativas de validación y de soporte asíncrono integrado. & \covCross \\ \bighline
\end{tabularx}%
}
\caption{Alternativas tecnológicas para el desarrollo del backend.}
\label{tab:comp_backend}
\end{table}

\subsubsection{Lenguaje del cliente web}
Se contrasta la seguridad del tipado frente a la velocidad de desarrollo directa en el frontend.

\begin{table}[H]
\centering
\small
\renewcommand{\arraystretch}{1.5}
\makebox[\textwidth][c]{%
\begin{tabularx}{1\textwidth}{l Y Y c}
\bighline
\textbf{Tecnología} & \textbf{Ventajas} & \textbf{Desventajas} & \textbf{Elección} \\ \bighline
\textbf{TypeScript} & Tipado estático estricto; mitiga errores con contratos claros para \textit{WebSockets}. & Requiere tiempo adicional en el ciclo de desarrollo para configurar interfaces. & \covCheck \\ \hline
JavaScript \textit{vanilla} & Curva de aprendizaje nula y compatibilidad nativa directa en el navegador. & Falta de seguridad de tipos en la mensajería asíncrona en tiempo real. & \covCross \\ \bighline
\end{tabularx}%
}
\caption{Alternativas para el lenguaje de programación frontend.}
\label{tab:comp_lenguaje}
\end{table}

\subsubsection{Frontend (interfaz de usuario)}
Se analiza la arquitectura de las interfaces y la velocidad en la compilación del entorno de desarrollo.

\begin{table}[H]
\centering
\small
\renewcommand{\arraystretch}{1.4}
\makebox[\textwidth][c]{%
\begin{tabularx}{1\textwidth}{l P{5.5cm} Y c}
\bighline
\textbf{Tecnología} & \textbf{Ventajas} & \textbf{Desventajas} & \textbf{Elección} \\ \bighline
\textbf{React + Vite} & \textit{Virtual DOM} reactivo siguiendo las pautas de React de Porcello y Banks \cite{porcello2020learning}, patrones de interfaz según MacDonald \cite{macdonald2019practical} y compilación \textit{HMR} instantánea. & Curva de aprendizaje inicial superior frente a opciones más básicas. & \covCheck \\ \hline
Angular & \textit{Framework} corporativo muy sólido, estructurado y altamente escalable. & Curva de aprendizaje muy empinada y verbosidad excesiva para el alcance del TFG. & \covCross \\ \hline
Vue.js & Alternativa equilibrada con un sistema de reactividad sencillo y progresivo. & Ecosistema y penetración en el mercado laboral inferior en comparación con React. & \covCross \\ \hline
Svelte & Máxima ligereza al eliminar el \textit{virtual DOM} y compilar a código nativo. & Comunidad reducida y menor disponibilidad de librerías de terceros maduras. & \covCross \\ \hline
Webpack & Estándar histórico para el empaquetado de aplicaciones web. & Procesos de construcción muy lentos en desarrollo local comparado con Vite. & \covCross \\ \bighline
\end{tabularx}%
}
\caption{Alternativas tecnológicas para la interfaz de usuario.}
\label{tab:comp_frontend}
\end{table}

\subsubsection{Estilización visual de la interfaz}
Se evalúa el equilibrio entre la velocidad de diseño y la sobrecarga estética en la aplicación del cliente.

\begin{table}[H]
\centering
\small
\renewcommand{\arraystretch}{1.4}
\makebox[\textwidth][c]{%
\begin{tabularx}{1\textwidth}{l Y Y c}
\bighline
\textbf{Tecnología} & \textbf{Ventajas} & \textbf{Desventajas} & \textbf{Elección} \\ \bighline
\textbf{CSS \textit{vanilla}} & Control total de layouts y animaciones sin dependencias ni sobrecarga. & Requiere la escritura manual de todos los estilos desde cero. & \covCheck \\ \hline
Tailwind CSS & Maquetación veloz mediante clases de utilidad atómicas directas. & Saturación excesiva del marcado JSX, perjudicando la legibilidad del código. & \covCross \\ \hline
Bootstrap & Componentes preconstruidos adaptables y fáciles de implementar. & Estética genérica y obsoleta; complejiza la personalización visual premium. & \covCross \\ \bighline
\end{tabularx}%
}
\caption{Alternativas tecnológicas para el diseño de estilos visuales.}
\label{tab:comp_estilos}
\end{table}

\subsubsection{Infraestructura y despliegue}
Se analizan los entornos de alojamiento priorizando la compatibilidad con canales bidireccionales y la viabilidad económica.

\begin{table}[H]
\centering
\small
\renewcommand{\arraystretch}{1.4}
\makebox[\textwidth][c]{%
\begin{tabularx}{1\textwidth + 0.3cm}{l P{5.4cm} Y c}
\bighline
\textbf{Tecnología} & \textbf{Ventajas} & \textbf{Desventajas} & \textbf{Elección} \\ \bighline
\textbf{Vercel + Render} & Plan gratuito sin tarjeta, GitOps automatizado y soporte ASGI para \textit{WebSockets}. & El plan gratuito de Render sufre de retardos iniciales por \textit{cold start}. & \covCheck \\ \hline
AWS & Control absoluto e ilimitado de la infraestructura de red. & Configuración compleja y riesgo de costes económicos imprevistos. & \covCross \\ \hline
Heroku & Entorno \textit{PaaS} históricamente óptimo para servidores de sockets interactivos. & Definición obsoleta por la eliminación por completo de sus planes gratuitos. & \covCross \\ \bighline
\end{tabularx}%
}
\caption{Alternativas para la infraestructura y el despliegue.}
\label{tab:comp_infraestructura}
\end{table}

\subsection{Limitaciones tecnológicas}\label{subsec:limitaciones_tecnologicas}

A pesar de la idoneidad del stack seleccionado, se han identificado las siguientes limitaciones técnicas inherentes a las decisiones tomadas:
\begin{itemize}
    \item \textbf{Cold start en Render (plan gratuito):} El contenedor del backend se desactiva tras 15 minutos sin recibir peticiones. Esto genera un retardo inicial de entre 30 y 50 segundos en la primera interacción de un usuario mientras el contenedor vuelve a iniciarse.
    \item \textbf{Concurrencia de escritura en SQLite:} SQLite bloquea la base de datos a nivel de archivo completo durante las operaciones de escritura. Aunque es imperceptible para un grupo estándar de alumnos, limitaría el escalado para miles de usuarios simultáneos registrando respuestas al unísono.
    \item \textbf{Falta de SSR (\textit{server-side rendering}) en el frontend:} Al ser una SPA de React pura, el indexado por motores de búsqueda (SEO) es limitado. Sin embargo, al tratarse de una aplicación académica de acceso cerrado mediante credenciales, esta limitación no compromete su viabilidad.
\end{itemize}

\subsection{Herramientas de soporte al desarrollo y control de calidad}\label{subsec:soporte_desarrollo}

La robustez del sistema y el cumplimiento de los estándares de calidad se articulan mediante las siguientes herramientas de soporte:

\begin{itemize}
    \item \textbf{Control de versiones y flujo de trabajo:} Se utiliza \textbf{Git} como sistema de control local y \textbf{GitHub} como repositorio remoto. La gestión de ramas aplica el modelo de ramificación de Driessen \cite{driessen2010git} (\textit{Git Flow}), manteniendo una rama estable (\texttt{main}), una de integración (\texttt{develop}) y ramas de características (\texttt{feat/*}). Siguiendo las pautas de organización de Westby \cite{westby2015git}, cada funcionalidad se desarrolla de forma aislada en su propia rama antes de fusionarse directamente en \texttt{develop}.
    \item \textbf{Gestión de tareas y proyectos:} Para el seguimiento del cronograma se utiliza \textbf{GitHub Projects} estructurado como un tablero \textit{Kanban}. Esto garantiza la trazabilidad del trabajo al asociar los \textit{commits} y ramas de desarrollo con su correspondiente tarea (\textit{issue}) e hito (\textit{milestone}).
    \item \textbf{Estandarización de commits:} Para garantizar un historial de Git limpio, semántico y legible, se integra \textbf{Commitizen}. Esta herramienta obliga al autor a seguir la convención de \textit{Conventional Commits}, categorizando de forma estricta cada cambio (por ejemplo, \textit{feat}, \textit{fix}, \textit{docs} o \textit{refactor}), lo que facilita la auditoría del código y la futura generación automatizada de diarios de cambios (\textit{changelogs}).
    \item \textbf{Entorno de desarrollo integrado (IDE):} El entorno principal de trabajo es \textbf{Visual Studio Code}, seleccionado por su ligereza y su ecosistema de extensiones. Entre ellas, se ha configurado \textbf{LaTeX Workshop} para centralizar la composición y compilación local de la memoria académica mediante los motores de \textbf{MiKTeX} y \textbf{Strawberry Perl}.
    \item \textbf{Gestión de entornos y dependencias:} En el entorno de backend se ha adoptado \textbf{uv}, una herramienta de gestión de paquetes y entornos virtuales para Python extremadamente rápida escrita en Rust, que sustituye de forma eficiente a herramientas tradicionales como \textit{pip} y \textit{poetry}. Para el entorno frontend se utiliza \textbf{npm} como gestor de paquetes estándar de Node.js.
    \item \textbf{Generación y lectura de códigos QR:} Para dar soporte al sistema de validación de presencia física, en el frontend se integra la librería \textbf{qrcode.react} para renderizar de manera eficiente y reactiva el código QR dinámico en la pantalla del alumno al finalizar el test. Por la parte del docente, se utiliza \textbf{html5-qrcode} para acceder a la cámara del dispositivo móvil y procesar el escaneo en tiempo real con baja latencia.
    \item \textbf{Linter y formateo:} En el backend se utiliza \textbf{Ruff}, \textit{linter} y formateador ultrarrápido programado en Rust que optimiza el cumplimiento de PEP 8 y unifica tareas que antes requerían múltiples librerías independientes (como \textit{Flake8}, \textit{Black} e \textit{isort}). En el frontend se emplea \textbf{ESLint} con \textit{flat config} para controlar la calidad de TypeScript junto a \textbf{Prettier} para garantizar un formateo de código unificado y consistente.
    \item \textbf{Pre-commit hooks:} Con la librería \textbf{pre-commit}, se interceptan los \textit{commits} locales para garantizar que solo se sube código limpio que supere de forma automatizada las reglas de formato, tipado y análisis estático antes de salir del entorno local.
    \item \textbf{Pruebas unitarias y de integración:} Se utiliza \textbf{pytest} en el backend junto a \textbf{pytest-cov} para el cálculo de cobertura, y \textbf{Vitest} con \textbf{React Testing Library} en el frontend para evaluar de manera aislada el comportamiento de los componentes. Para flujos de usuario extremo a extremo (E2E), se integra \textbf{Playwright}, permitiendo simular múltiples navegadores concurrentes.
    \item \textbf{Pruebas de carga:} Se emplea \textbf{Locust} para verificar de manera empírica la capacidad del servidor y medir la latencia del protocolo \textit{WebSocket} al soportar la avalancha de respuestas concurrentes de los alumnos.
    \item \textbf{SonarCloud:} Se ha configurado este auditor externo de análisis estático como parte del flujo de integración y entrega continuas (CI/CD) en GitHub Actions, aplicando los principios de canalizaciones de despliegue automatizadas descritos por Farley \cite{farley2021continuous} para exigir una meta de calidad estricta (\textit{quality gate} A, duplicación $<$ 3\%, cobertura de tests $>$ 80\%).
\end{itemize}

\section{Conclusiones del estudio previo}\label{sec:conclusiones_estudio_previo}

El análisis del estado del arte y el estudio comparativo detallados en este capítulo ponen de manifiesto una brecha crítica en el ámbito de la evaluación gamificada en clase. A pesar de la madurez de las soluciones analizadas, el diagnóstico del mercado revela que ninguna alternativa resuelve con éxito el problema del \textbf{fraude por suplantación a distancia} ni ofrece flujos optimizados que mitiguen la carga administrativa del profesorado. La falta de mecanismos de \textbf{control de asistencia eficaces} en las herramientas actuales y la rigidez de los entornos virtuales tradicionales justifican plenamente la necesidad de desarrollar una solución integral que logre unificar \textbf{inmediatez, control de presencia real e interoperabilidad}.

Como respuesta directa a este escenario de oportunidad, la propuesta de valor de \textit{Quizzie} se consolida a través de los factores diferenciadores definidos. El núcleo innovador del proyecto se materializa en el diseño de un \textbf{flujo de validación mediante códigos QR}, devolviendo al docente el control físico del aula sin penalizar el ritmo de la sesión. Asimismo, para contrarrestar la monotonía organizativa identificada en la competencia, la plataforma complementa su arquitectura con herramientas secundarias destinadas a agilizar la gestión diaria: la \textbf{generación de cuestionarios asistida por Inteligencia Artificial} y la \textbf{exportación directa de calificaciones} en formatos universales (CSV y PDF). Con ello, el proyecto trasciende el enfoque puramente lúdico y se posiciona como una \textbf{herramienta formal y eficiente} para el ecosistema educativo.

Para hacer esto posible, se ha elegido a conciencia un stack tecnológico donde cada pieza responde directamente a los retos de rendimiento del proyecto. La combinación de una base de datos relacional y ligera en local junto a un \textbf{servidor backend asíncrono} permite gestionar la \textbf{concurrencia de los cuestionarios} en el aula con total fluidez, demostrando que es viable construir un sistema interactivo robusto sin depender de infraestructura costosa en la nube. Esta selección técnica se cierra de forma lógica con el \textbf{ecosistema de automatización integrado}; al coordinar el formateo, las pruebas de rendimiento y la auditoría continua antes de cada despliegue, se garantiza que \textit{Quizzie} no solo funcione bien en teoría, sino que sea un software \textbf{estable, mantenible y listo para su uso real} en clase.